\documentclass[/work/src/main.tex]{subfiles}

\begin{document}
\section{Personalizaciones de la Plantilla}
\label{sec:personalizaciones}

Esta sección documenta las capacidades que la plantilla incorpora sobre LaTeX estándar.

\subsection{Macros Personalizados}
\label{sec:macros}

El archivo \texttt{settings/macros.tex} define comandos reutilizables:

\begin{itemize}
    \item \texttt{\textbackslash marketterm\{texto\}} — marca un término en cursiva: \marketterm{ejemplo de término}.
\end{itemize}

\subsection{Entornos de Teoremas}
\label{sec:teoremas}

\emph{Definidos en \texttt{settings/theorems.tex}.}

\begin{theorem}
    Enunciado de un teorema numerado por sección.
\end{theorem}

\begin{lemma}
    Lema auxiliar que apoya la demostración de teoremas.
\end{lemma}

\begin{proposition}
    Proposición que puede ser demostrada.
\end{proposition}

\begin{example}
    Ejemplo ilustrativo de un concepto.
\end{example}

\subsection{Referencias Cruzadas Inteligentes}
\label{sec:cleveref}

El paquete \texttt{cleveref} (configurado en \texttt{settings/links.tex}) genera automáticamente el tipo de referencia. Por ejemplo:

\begin{itemize}
    \item \texttt{\textbackslash cref\{sec:macros\}} → \cref{sec:macros}
    \item \texttt{\textbackslash Cref\{sec:teoremas\}} → \Cref{sec:teoremas}
    \item \texttt{\textbackslash cref\{sec:decoraciones\}} → \cref{sec:decoraciones}
\end{itemize}

El idioma español se detecta automáticamente desde \texttt{babel}.

\subsection{Comillas Contextuales}
\label{sec:csquotes}

El paquete \texttt{csquotes} adapta las comillas al idioma del documento:

\begin{quote}
    \enquote{Esto es una cita con comillas inteligentes.}
\end{quote}

\subsection{Estilo de Secciones}
\label{sec:estilo}

\emph{Configurado en \texttt{settings/styles.tex}.}

Los títulos de sección utilizan colores personalizados mediante \texttt{titlesec}:

\begin{itemize}
    \item \textcolor{sectioncolor}{\textbf{Secciones}} en azul oscuro.
    \item \textcolor{subsectioncolor}{\textbf{Subsecciones}} en verde azulado.
    \item \textcolor{subsubsectioncolor}{\textbf{Subsubsecciones}} en púrpura.
\end{itemize}

\subsection{Encabezado y Pie de Página}
\label{sec:header}

\emph{Definido en \texttt{settings/header.tex}} mediante \texttt{fancyhdr}:

\begin{itemize}
    \item Cabecera izquierda: autor del documento.
    \item Cabecera derecha: nombre de la sección actual.
    \item Pie centrado: número de página.
\end{itemize}

\subsection{Líneas Decorativas}
\label{sec:decoraciones}

\emph{Definidas en \texttt{settings/decorations.tex}} con \texttt{tikz}.

\begin{itemize}
    \item \texttt{\textbackslash andreasline} — línea decorativa gruesa con opacidad variable.
    \item \texttt{\textbackslash andreaslineslim} — versión delgada para espacios reducidos.
\end{itemize}

Ejemplo:

\andreasline

\subsection{Mejoras Tipográficas}
\label{sec:microtype}

El paquete \texttt{microtype} se aplica automáticamente a todo el documento, mejorando el espaciado entre caracteres y la justificación del texto sin intervención manual.

\end{document}
